\documentclass[10pt]{article}

% Specify the margins and text size.
\setlength{\textwidth}{6.5in}
\setlength{\textheight}{9.5in}
\setlength{\oddsidemargin}{0pt}
\setlength{\evensidemargin}{0pt}
\setlength{\topmargin}{0pt}
\setlength{\hoffset}{-.05in}
\setlength{\voffset}{-.8in}

\setlength{\parskip}{5pt}
\setlength{\parindent}{0pt}

% Load some fonts and symbol packages
\usepackage{latexsym}
\usepackage{pifont}       % contains 'star' symbol for counterinsurgency handbook title
\usepackage{yfonts} 
\usepackage{amsmath}
\usepackage{amsfonts}

\usepackage{graphicx}     % actually, this is loaded by pstricks
\usepackage[T1]{fontenc}
\usepackage{ifthen}
\usepackage{pstricks,pst-grad,pst-text,pst-node,multido,pst-plot,calc,pst-3dplot}
%\usepackage[all]{xy}
%\usepackage{animate}

% The hyperref package inserts links into the pdf.
\definecolor{MyLinkColor}{rgb}{.1,.2,1}
\definecolor{MyCiteColor}{rgb}{.1,1,.2}
\definecolor{MyURLColor}{rgb}{.4,.4,.4}
\usepackage[backref=true,pagebackref=false,hyperindex,colorlinks=true,
  linkcolor=MyLinkColor,urlcolor=MyURLColor]{hyperref}


% The tweaklist package is something I found on the web.  It provides a simple interface
% for making changes to spacing used in the itemize and enumerate environments.  Comment
% this out if you don't care to use tweaklist.
\usepackage{tweaklist}
\renewcommand{\itemhook}{\setlength{\parskip}{2pt}\setlength{\parsep}%
{1pt}\setlength{\topsep}{0pt}\setlength{\itemsep}{0pt}}

\newcommand{\U}{\underline{\hspace{5pt}}}

\usepackage{listings}
\newcounter{EX}\setcounter{EX}{1}
\newcommand{\EXERCISE}{\arabic{EX}.\stepcounter{EX} }

\begin{document}
\pagestyle{empty}
\lstset{language=R, showspaces=false, showstringspaces=false}

\href{http://www.su.edu}{\includegraphics[height=1.75cm]{sulogo.eps}}
\vspace{-1.69cm}

{\small \hfill
\begin{tabular}{cl}
& Math 207\\ & Introduction to  Statistics\\
\end{tabular}
}
\setlength{\baselineskip}{1.05\baselineskip}

\begin{center}
\textbf{\large  Basic Histogram Exercises}
\end{center}

1. The histogram below shows the distribution of final scores in a certain
class.

\hspace{10pt} a) Which block represents the people who scored between
  60 and 80?

\hspace{10pt} b) Ten percent scored between 20 and 40.  About what percentage
  scored between 40 and 60?

\hspace{10pt} c) About what percentage scored over 60?

\begin{center}
\begin{pspicture}(-1,-0.5)(5,3)
\psset{yunit=0.8}
\psframe(0,0)(1,1)\rput(0.5,0.5){D}
\psframe(1,0)(2,2)\rput(1.5,1.0){C}
\psframe(2,0)(3,4)\rput(2.5,2.5){B}
\psframe(3,0)(4,3)\rput(3.5,1.5){A}
\rput(0,-0.3){\small 20}
\rput(1,-0.3){\small 40}
\rput(2,-0.3){\small 60}
\rput(3,-0.3){\small 80}
\rput(4,-0.3){\small 100}
\rput(2,-1){\small Final score}
\psset{linewidth=0.02}
\psline(-0.5,0)(-0.5,4)
\psline(-0.5,0)(-0.3,0)\rput(-0.8,0){\small 0}
\psline(-0.5,1)(-0.3,1)\rput(-0.8,1){\small 1}
\psline(-0.5,2)(-0.3,2)\rput(-0.8,2){\small 2}
\psline(-0.5,3)(-0.3,3)\rput(-0.8,3){\small 3}
\psline(-0.5,4)(-0.3,4)\rput(-0.8,4){\small 4}
\end{pspicture}
\end{center}

2. Below are sketches of histograms for test scores in three different classes.  
 Scores range from 0 to 100; a passing score was 50.  
For each class, was the percent who  passed about 50\%, well over 50\%, or
well under~50\%?

\begin{center}
{\setlength{\tabcolsep}{20pt}\begin{tabular}{ccc}
\begin{pspicture}(0,0)(4,1)
\psset{xunit=4, yunit=3}
\psline(0,0)(.6,.4)(.7,.48)(.75,.5)(.8,.48)(1,0)
\psset{linewidth=0.02}
\psline(0,0)(1,0)
\psline(0.25,0.05)(0.25,-0.05)\rput(0.25,-0.15){\small 25}
\psline(0.5 ,0.05)(0.5 ,-0.05)\rput(0.5 ,-0.15){\small 50}
\psline(0.75,0.05)(0.75,-0.05)\rput(0.75,-0.15){\small 75}
\psline(1   ,0.05)(1   ,-0.05)\rput(1   ,-0.15){\small 100}
\rput(0,0.25){(a)}
\end{pspicture}
&
\begin{pspicture}(0,0)(4,1)
\psset{xunit=4, yunit=3}
\psline(0,0)(0.1,0.38)(0.12,0.4)(0.38,0.4)(0.4,0.38)(0.5,0.05)(0.9,0.05)(0.93,0.5)(0.95,0.52)(0.97,0.5)(1,0)
\psset{linewidth=0.02}
\psline(0,0)(1,0)
\psline(0.25,0.05)(0.25,-0.05)\rput(0.25,-0.15){\small 25}
\psline(0.5 ,0.05)(0.5 ,-0.05)\rput(0.5 ,-0.15){\small 50}
\psline(0.75,0.05)(0.75,-0.05)\rput(0.75,-0.15){\small 75}
\psline(1   ,0.05)(1   ,-0.05)\rput(1   ,-0.15){\small 100}
\rput(-0.05,0.25){(b)}
\end{pspicture}
&
\begin{pspicture}(0,0)(4,1)
\psset{xunit=4, yunit=3}
\psline(0,0)(0.12,0.05)(0.25,0.15)(0.45,0.48)(0.5,0.5)(0.55,0.48)(0.75,0.15)(0.88,0.05)(1,0)
\psset{linewidth=0.02}
\psline(0,0)(1,0)
\psline(0.25,0.05)(0.25,-0.05)\rput(0.25,-0.15){\small 25}
\psline(0.5 ,0.05)(0.5 ,-0.05)\rput(0.5 ,-0.15){\small 50}
\psline(0.75,0.05)(0.75,-0.05)\rput(0.75,-0.15){\small 75}
\psline(1   ,0.05)(1   ,-0.05)\rput(1   ,-0.15){\small 100}
\rput(0,0.25){(c)}
\end{pspicture}
\end{tabular}}
\end{center}
\vspace{.5in}

3. One class in exercise 2 had two quite distinct groups of students, with one 
group doing rather poorly on the test, and the other group doing very well.  
What class was it?
\vspace{.5in}

4. In class (b) of exercise 2, were there more people with scores in the range
40--50 or 90--100?
\vspace{.5in}


5. Three people plot histograms for the weights of subjects in a study, 
using the density scale.  Only one is right.  Which one, and why?

\begin{center}
\begin{tabular}{ccccc}
\begin{pspicture}(0,0)(4,1.5)
\psline(0,0.2)(4,0.2)
\psline(0,0.1)(0,0.3) \rput[t](0,-0.1){100}
\psline(1,0.1)(1,0.3) 
\psline(2,0.1)(2,0.3) \rput[t](2,-0.1){150}
\psline(3,0.1)(3,0.3) 
\psline(4,0.1)(4,0.3) \rput[t](4,-0.1){200}
\psarc[linecolor=blue,linewidth=0.02](2,-0.3){2.06}{14}{166}
\rput(2,0.85){(i)}
\end{pspicture}
&\hspace{.4in} &
\begin{pspicture}(0,0)(4,1.5)
\psline(0,0.2)(4,0.2)
\psline(0,0.1)(0,0.3) \rput[t](0,-0.1){100}
\psline(1,0.1)(1,0.3) 
\psline(2,0.1)(2,0.3) \rput[t](2,-0.1){150}
\psline(3,0.1)(3,0.3) 
\psline(4,0.1)(4,0.3) \rput[t](4,-0.1){200}
\psarc[linecolor=blue,linewidth=0.02](2,-0.3){2.06}{14}{166}
\rput[t](2,-0.6){Weight (lbs)}
\psline(0,0.2)(0,1.5)
\psline(-0.1,0.2)(0.1,0.2)\rput[r](-0.2,0.2){0}
\psline(-0.1,0.85)(0.1,0.85)\rput[r](-0.2,0.85){1}
\psline(-0.1,1.5)(0.1,1.5)\rput[r](-0.2,1.5){2}
\rput{90}(-0.75,0.85){\% per lb}
\rput(2,0.85){(ii)}
\end{pspicture}
&\hspace{.4in} &
\begin{pspicture}(0,0)(4,1.5)
\psline(0,0.2)(4,0.2)
\psline(0,0.1)(0,0.3) \rput[t](0,-0.1){100}
\psline(1,0.1)(1,0.3) 
\psline(2,0.1)(2,0.3) \rput[t](2,-0.1){150}
\psline(3,0.1)(3,0.3) 
\psline(4,0.1)(4,0.3) \rput[t](4,-0.1){200}
\psarc[linecolor=blue,linewidth=0.02](2,-0.3){2.06}{14}{166}
\rput[t](2,-0.6){Weight (lbs)}
\psline(0,0.2)(0,1.5)
\psline(-0.1,0.2)(0.1,0.2)\rput[r](-0.2,0.2){0}
\psline(-0.1,0.85)(0.1,0.85)\rput[r](-0.2,0.85){1}
\psline(-0.1,1.5)(0.1,1.5)\rput[r](-0.2,1.5){2}
\rput{90}(-0.75,0.85){lb per \%}
\rput(2,0.85){(iii)}
\end{pspicture}
\end{tabular}
\end{center}
\bigskip


\vfill
\eject
6. Two histograms are sketched below.  One shows the distribution of age at 
death from natural causes (heart disease, cancer, and so forth). The other shows
age at death from trauma (accident, murder, suicide).  Which is which and why?

\begin{center}
{\setlength{\tabcolsep}{20pt}\begin{tabular}{cc}
\begin{pspicture}(0,0)(4,1)
\psset{xunit=4, yunit=3}
\psline(0,0)(.5,.13)(.55,.16)(.6,.2)(.65,.24)(.7,.3)(.75,.35)(.8,.38)(.85,.35)%
(.9,.3)(.95,.22)(1,0)
\psset{linewidth=0.02}
\psline(0,0)(1,0)
\psline(0.25,0.05)(0.25,-0.05)\rput(0.25,-0.15){\small 25}
\psline(0.5 ,0.05)(0.5 ,-0.05)\rput(0.5 ,-0.15){\small 50}
\psline(0.75,0.05)(0.75,-0.05)\rput(0.75,-0.15){\small 75}
\psline(1   ,0.05)(1   ,-0.05)\rput(1   ,-0.15){\small 100}
\rput(0,0.25){(a)}
\end{pspicture}
&
\begin{pspicture}(0,0)(4,1)
\psset{xunit=4, yunit=3}
\psline(1,0)(.5,.13)(.45,.16)(.4,.2)(.35,.24)(.3,.3)(.25,.35)(.2,.38)(.15,.35)%
(.1,.3)(.05,.22)(0,0)
\psset{linewidth=0.02}
\psline(0,0)(1,0)
\psline(0.25,0.05)(0.25,-0.05)\rput(0.25,-0.15){\small 25}
\psline(0.5 ,0.05)(0.5 ,-0.05)\rput(0.5 ,-0.15){\small 50}
\psline(0.75,0.05)(0.75,-0.05)\rput(0.75,-0.15){\small 75}
\psline(1   ,0.05)(1   ,-0.05)\rput(1   ,-0.15){\small 100}
\rput(-0.05,0.25){(b)}
\end{pspicture}
\end{tabular}}
\end{center}
\vspace{1.5in}


7. A histogram of monthly wages for part-time employees is shown 
below (densities are marked
in parentheses).  Nobody earned more than \$1,000 per month.  

\begin{center}
\begin{pspicture}(0,-0.8)(10,2.3)
\psset{yunit=0.1}
\psframe(0,0)(1,10)\rput[b](0.5,10.1){(10)}
\psframe(1,0)(2,20)\rput[b](1.5,20.1){(20)}
\psframe(5,0)(10,5)\rput[b](7.5,5.1){(5)}
\psaxes[Dy=10](0,0)(10,20)
\rput{90}(-1.1,10){\% per \$100}
\rput(5,-10){Wages (hundreds of dollars)}
\end{pspicture}
\end{center}

\hspace{10pt} a) The block over the class interval from \$200 to \$500 is missing.
How tall must it be?
\vspace{.4in}

\hspace{10pt} b) What percent of the employees made at least \$200 per month?
\vspace{.4in}

\hspace{10pt} c) About what percent of the employees made between \$600 and
\$800 per month?
\vfill
\eject

\begin{center}
\textbf{\large Solutions}
\end{center}
\bigskip

1.

(a) Block B represents the people who scored between 60 and 80.

(b) Since 10\% scored between 20 and 40, the area of block D is 10\%.  
The area of block C represents the percent of people who scored between 40 and 60.
It looks like this block is about twice as tall as block D.  Both block have the 
same width.  So the answer is 20\%.

(c) Blocks D and C make up a total of 30\%.  The percentage over 60 is the other two blocks (B and A).
So, their area is $100\% - 30\% = 70\%$.
\bigskip

2.

(a) The percent who passed was well over 50\% since most of the area is above the score of 50.

(b) The percent who passed was well under 50\% since most of the area is below 
the score of 50.

(c) The percent who passed was around 50\% since the area is symmetric around 
the score of 50.
\bigskip

3. Class (b) had two distinct groups.  One small group did very well.  
A much larger group (larger area)
failed because their scores were below 50.
\bigskip

4. Think of the rectangle for the group that scored 90--100 and the rectangle for the group that scored 40--50.
Both rectangles have width 10.  The 90--100 rectangle is taller so it has a bigger area.  So, more students
scored 90--100 than scored 40--50.
\bigskip

5. Histogram (ii) is correct.  It is the only one with the correct density scale on the vertical axis.
Histogram (i) doesn't have a scale on either axis.  On histogram (iii), the units on the vertical
scale are wrong.  The correct units are \% per lb  instead of lb per \%.
\bigskip

6.
Histogram (a) shows the distribution of ages from natural causes. Few young people will die of such causes
so there is little area on the left side. Histogram (b) shows the distribution of ages from trauma. This is a
more likely cause of death for young people than for old people so the histogram has a lot of area on the
left side but not so much on the right.
\bigskip

7. 

a) The areas of the three given blocks are 10\%, 20\% and 25\%.  This gives a 
total of 55\%.  The missing box must have an area of  45\%.  Since its width
is 3, its height must be 15.  

b) 100\% - (10\% + 20\%) = 100\% - 30\% = 70\%.  Alternatively, 45\% + 25\% = 70\%.

c)  $w\times h = 2\times 5 = 10\%$.
\vfill
\eject

\end{document}
